\lecture{13}{Инерционные методы}

\sourcevideo
    {https://www.youtube.com/playlist?list=PLOzRYVm0a65fhKDS8cYIC7YCuVGJ4FOJx}
    {Week 4: Lecture 13: Accelerate the Convergence Even Further}

Для выпуклой \(L\)-гладкой функции градиентный спуск имеет скорость
\(\mathcal O(1/k)\), а при PL-условии --- линейную скорость с
зависимостью от числа обусловленности \(\kappa=L/\mu\). В этой лекции
ускорение достигается не усилением предположений о функции, а
изменением самого алгоритма: к градиентному шагу добавляется память о
предыдущем движении.

\section{Мотивация импульса и исходные оценки}

Рассмотрим задачу
\[
    \min_{x\in\RR^n}f(x)
\]
и градиентный спуск
\[
    x^{k+1}=x^k-\frac1L\nabla f(x^k).
\]
Для выпуклой \(L\)-гладкой функции
\[
    f(x^k)-f^\star
    =
    \mathcal O\!\left(\frac1k\right),
\]
поэтому для достижения ошибки не более \(\varepsilon\) требуется
\[
    k
    =
    \mathcal O\!\left(\frac1\varepsilon\right).
\]
Если функция удовлетворяет PL-неравенству
\[
    \frac12\norm{\nabla f(x)}^2
    \ge
    \mu\bigl(f(x)-f^\star\bigr),
\]
то
\[
    f(x^k)-f^\star
    \le
    \left(1-\frac\mu L\right)^k
    \bigl(f(x^0)-f^\star\bigr).
\]
В сильно выпуклом случае
\[
    \kappa=\frac L\mu,
\]
и сложность равна
\[
    \mathcal O\!\left(
        \kappa\log\frac1\varepsilon
    \right).
\]

Причина замедления на плохо обусловленной функции видна по геометрии
линий уровня. Градиент ортогонален линии уровня, поэтому в вытянутой
долине последовательные шаги могут многократно пересекать её поперёк,
медленно продвигаясь к минимуму. Обычный градиентный спуск использует
только \(x^k\) и \(\nabla f(x^k)\); направление предыдущего
перемещения полностью теряется. Инерционный метод сохраняет часть
вектора
\[
    x^k-x^{k-1}
\]
и тем самым может ослаблять поперечные колебания, не уничтожая
накопленное движение вдоль долины.

\section{Метод тяжёлого шара}

Метод тяжёлого шара Поляка задаётся правилом
\[
    x^{k+1}
    =
    x^k
    -
    \alpha\nabla f(x^k)
    +
    \beta(x^k-x^{k-1}),
\]
где
\[
    \alpha>0,
    \qquad
    \beta\ge0.
\]
При \(\beta=0\) получается обычный градиентный спуск. Если ввести
скорость
\[
    v^k=x^k-x^{k-1},
\]
то алгоритм принимает вид
\[
    v^{k+1}
    =
    \beta v^k-\alpha\nabla f(x^k),
    \qquad
    x^{k+1}=x^k+v^{k+1}.
\]

\begin{algorithm}[H]
\caption{Метод тяжёлого шара}
\label{alg:heavy-ball}
\begin{algorithmic}[1]

\Require
\(x^0\in\RR^n\), \(v^0=0\),
параметры \(\alpha>0\), \(\beta\ge0\)

\For{\(k=0,1,2,\ldots\)}

    \State
    \(g^k\gets\nabla f(x^k)\)

    \State
    \(v^{k+1}\gets\beta v^k-\alpha g^k\)

    \State
    \(x^{k+1}\gets x^k+v^{k+1}\)

\EndFor

\end{algorithmic}
\end{algorithm}

Метод использует два последовательных состояния, но остаётся методом
первого порядка: Гессиан не вычисляется. Одна итерация требует одного
градиента, как и обычный градиентный спуск, однако дополнительно нужно
хранить один вектор размерности \(n\).

В лекции при надлежащем выборе \(\alpha\) и \(\beta\) приводятся
оценки
\[
    \mathcal O\!\left(\frac1\varepsilon\right)
\]
для общего выпуклого гладкого случая и
\[
    \mathcal O\!\left(
        \sqrt{\frac L\mu}
        \log\frac1\varepsilon
    \right)
    =
    \mathcal O\!\left(
        \sqrt\kappa
        \log\frac1\varepsilon
    \right)
\]
для сильно выпуклого случая. Подробное доказательство и точная область
допустимых параметров в лекции не приводятся; наиболее прозрачный
анализ метода тяжёлого шара получается для квадратичных функций.

\section{Метод Нестерова}

Метод Нестерова также использует импульс, но градиент вычисляется в
экстраполированной точке. Сначала строится
\[
    y^k
    =
    x^k+\beta_k(x^k-x^{k-1}),
\]
после чего выполняется шаг
\[
    x^{k+1}
    =
    y^k-\frac1L\nabla f(y^k).
\]
У тяжёлого шара градиент берётся в \(x^k\), а у метода Нестерова ---
в прогнозной точке \(y^k\).

Если \(f\) является \(L\)-гладкой и \(\mu\)-сильно выпуклой, можно
взять постоянный коэффициент
\[
    \beta
    =
    \frac{\sqrt L-\sqrt\mu}
         {\sqrt L+\sqrt\mu}
    =
    \frac{\sqrt\kappa-1}
         {\sqrt\kappa+1}.
\]
Тогда
\[
    y^k
    =
    x^k+\beta(x^k-x^{k-1}),
\]
\[
    x^{k+1}
    =
    y^k-\frac1L\nabla f(y^k),
\]
и для подходящей константы \(C>0\)
\[
    f(x^k)-f^\star
    \le
    C
    \left(
        1-\sqrt{\frac\mu L}
    \right)^k.
\]
Следовательно,
\[
    k
    =
    \mathcal O\!\left(
        \sqrt{\frac L\mu}
        \log\frac1\varepsilon
    \right).
\]

Удобно использовать стандартную эквивалентную параметризацию через \(t_k\).

\begin{algorithm}[H]
\caption{Ускоренный градиентный метод для выпуклой функции}
\label{alg:nesterov-convex}
\begin{algorithmic}[1]

\Require
\(x^{-1}=x^0\), константа \(L>0\), \(t_0=1\)

\For{\(k=0,1,2,\ldots\)}

    \If{\(k=0\)}

        \State
        \(y^0\gets x^0\)

    \Else

        \State
        \[
            y^k
            \gets
            x^k
            +
            \frac{t_{k-1}-1}{t_k}
            (x^k-x^{k-1})
        \]

    \EndIf

    \State
    \(x^{k+1}
        \gets
        y^k-\dfrac1L\nabla f(y^k)\)

    \State
    \[
        t_{k+1}
        \gets
        \frac{1+\sqrt{1+4t_k^2}}2
    \]

\EndFor

\end{algorithmic}
\end{algorithm}

\section{Скорости и оптимальность}

Для ускоренного метода Нестерова в выпуклом случае выполняется оценка
\[
    f(x^k)-f^\star
    \le
    \frac{
        2L\norm{x^0-x^\star}^2
    }{
        (k+1)^2
    }.
\]
Поэтому
\[
    f(x^k)-f^\star
    =
    \mathcal O\!\left(\frac1{k^2}\right),
\]
а для достижения точности \(\varepsilon\) достаточно
\[
    k
    =
    \mathcal O\!\left(
        \sqrt{
            \frac{
                L\norm{x^0-x^\star}^2
            }{
                \varepsilon
            }
        }
    \right).
\]
При фиксированных \(L\) и начальном расстоянии это даёт зависимость
\[
    \mathcal O\!\left(
        \frac1{\sqrt\varepsilon}
    \right)
\]
вместо
\[
    \mathcal O\!\left(
        \frac1\varepsilon
    \right)
\]
у градиентного спуска.

Основные оценки лекции имеют вид
\[
\begin{array}{c|c|c}
    \text{метод}
    &
    \text{выпуклая функция}
    &
    \text{сильно выпуклая функция}
    \\
    \hline
    \text{градиентный спуск}
    &
    \mathcal O(1/\varepsilon)
    &
    \mathcal O\!\left(
        \kappa\log(1/\varepsilon)
    \right)
    \\
    \text{тяжёлый шар}
    &
    \mathcal O(1/\varepsilon)
    &
    \mathcal O\!\left(
        \sqrt\kappa\log(1/\varepsilon)
    \right)
    \\
    \text{метод Нестерова}
    &
    \mathcal O(1/\sqrt\varepsilon)
    &
    \mathcal O\!\left(
        \sqrt\kappa\log(1/\varepsilon)
    \right)
\end{array}
\]

В модели оракула первого порядка существуют нижние оценки
\[
    \Omega\!\left(
        \sqrt{
            \frac{
                L\norm{x^0-x^\star}^2
            }{
                \varepsilon
            }
        }
    \right)
\]
для выпуклых \(L\)-гладких функций и
\[
    \Omega\!\left(
        \sqrt{\frac L\mu}
        \log\frac1\varepsilon
    \right)
\]
для \(L\)-гладких \(\mu\)-сильно выпуклых функций. Метод Нестерова
достигает этих порядков. Оптимальность относится именно к данной
модели доступа через значения функции и градиенты; дополнительная
структура задачи или использование информации второго порядка могут
изменить нижние оценки.

\section{Переход к непрерывному времени}

Градиентный спуск с шагом \(\eta>0\) можно записать как
\[
    \frac{x^{k+1}-x^k}{\eta}
    =
    -\nabla f(x^k).
\]
При формальном предельном переходе \(\eta\to0\) получается градиентный
поток
\[
    \dot x(t)
    =
    -\nabla f(x(t)).
\]

\begin{definition}[Градиентный поток]
Динамическая система
\[
    \dot x=-\nabla f(x)
\]
называется градиентным потоком функции \(f\).
\end{definition}

Для
\[
    f(x)=\frac12x^2
\]
имеем
\[
    \dot x=-x
\]
и
\[
    x(t)=e^{-t}x(0).
\]
Более общая система
\[
    \dot x=-ax,
    \qquad
    a>0,
\]
имеет решение
\[
    x(t)=e^{-at}x(0).
\]
В непрерывной модели увеличение \(a\) формально ускоряет затухание.
Однако дискретизация
\[
    x^{k+1}
    =
    (1-\eta a)x^k
\]
устойчива только при
\[
    0<\eta a<2.
\]
Следовательно, простое масштабирование непрерывного поля не даёт
бесплатного ускорения дискретного метода: рост коэффициента требует
уменьшения шага.

Глобальная \(L\)-гладкость особенно важна для анализа конечного шага,
поскольку она контролирует изменение градиента. В непрерывном анализе
некоторые оценки устойчивости можно получить без глобальной
константы \(L\), но условия регулярности поля всё равно нужны для
существования и единственности классического решения.

\section{Нормированный поток и практические ограничения}

Рассмотрим нормированный поток
\[
    \dot x
    =
    -
    \frac{
        \nabla f(x)
    }{
        \norm{\nabla f(x)}
    },
    \qquad
    \nabla f(x)\ne0.
\]
Его скорость удовлетворяет
\[
    \norm{\dot x}=1
\]
и не уменьшается около стационарной точки. Для
\[
    f(x)=\frac12x^2
\]
получаем
\[
    \dot x=-\operatorname{sign}(x),
\]
а решение до достижения нуля имеет вид
\[
    x(t)
    =
    \operatorname{sign}(x(0))
    \max\{
        \abs{x(0)}-t,
        0
    \}.
\]
Равновесие достигается за конечное время
\[
    T=\abs{x(0)}.
\]

В точке \(\nabla f(x)=0\) нормированное поле не определено, поэтому
нужно отдельно задать продолжение динамики или использовать
дифференциальное включение. После дискретизации возможны колебания
около минимума.

Все три дискретных метода используют один полный градиент за
итерацию. Градиентный спуск не требует дополнительного инерционного
состояния, а тяжёлый шар и метод Нестерова хранят ещё один вектор.
Ускорение по числу итераций не всегда совпадает с ускорением по
времени и памяти.

В распределённой задаче один градиентный шаг может включать
коммуникационный раунд, поэтому уменьшение числа итераций может
непосредственно сокращать число обменов, если стоимость одного шага
существенно не возрастает.

Непрерывная интерпретация связывает стационарные точки функции с
равновесиями динамической системы, сходимость алгоритма ---
с устойчивостью, а скорость оптимизации --- со скоростью затухания.
Следующая лекция вводит функции Ляпунова и формализует этот переход.
